\section{Conclusiones}

El presente trabajo de grado tuvo como objetivo general desarrollar un prototipo web que apoyara la comprensión de los contenidos de la asignatura Fundamentos de Interpretación y Compilación de Lenguajes de Programación (FLP) de la Universidad del Valle, sede Tuluá, mediante la visualización de estructuras internas y el uso de lenguajes orientados por sintaxis. A partir del diagnóstico de dificultades conceptuales, el diseño y la implementación del prototipo FLP Viewer, y su posterior evaluación con estudiantes de la asignatura, se considera que este objetivo fue alcanzado.

\subsection{Cumplimiento de los objetivos específicos}

El primer objetivo específico, orientado a identificar las dificultades conceptuales que enfrentan los estudiantes de la asignatura, se desarrolló en el Capítulo~\ref{cap3} mediante la aplicación de encuestas a 36 estudiantes y a los docentes responsables de los distintos niveles de la asignatura. El diagnóstico permitió confirmar que las dificultades no se concentran en la comprensión teórica general, sino específicamente en el segundo resultado de aprendizaje del microcurrículo, relacionado con ambientes de ejecución, estado y ligadura de variables, donde el 77.8\,\% de los estudiantes reportó dificultades explícitas. Este hallazgo orientó directamente la selección de los cinco indicadores de logro que sirvieron como base conceptual para el prototipo.

El segundo objetivo, referido a la caracterización de los fundamentos teóricos y técnicos sobre gramáticas formales, análisis sintáctico y árboles de sintaxis abstracta, se desarrolló en el Capítulo~\ref{cap2} mediante la construcción de un marco teórico que integró los fundamentos técnicos de la interpretación de lenguajes, las teorías del aprendizaje aplicadas a la enseñanza de compiladores y los principios de usabilidad en herramientas educativas interactivas. La búsqueda sistemática de literatura documentada en el Anexo~\ref{annex:slr} permitió, adicionalmente, identificar los vacíos que fundamentaron el aporte específico de este trabajo frente a las herramientas existentes.

El tercer objetivo, relacionado con el diseño de la arquitectura web y las interfaces de usuario, se desarrolló en el Capítulo~\ref{cap4}, donde se especificaron los requisitos funcionales y no funcionales del sistema y se definieron las decisiones de arquitectura, las tecnologías empleadas y las estrategias de visualización orientadas a mantener en un único entorno las representaciones que el estudiante debe relacionar para comprender el proceso de interpretación.

El cuarto objetivo, correspondiente a la implementación de un prototipo funcional, se desarrolló en el Capítulo~\ref{cap5}. El sistema implementado permite al estudiante definir gramáticas en notación BNF, generar y visualizar de forma interactiva el AST resultante, y observar los cambios en el ambiente de ejecución durante la evaluación paso a paso de un programa, cumpliendo con los ocho requisitos funcionales especificados en el Anexo~\ref{annex:rf}.

El quinto objetivo, relacionado con la evaluación de la usabilidad y funcionalidad del prototipo mediante pruebas con estudiantes y el docente responsable, se desarrolló en la Sección~\ref{sec:evaluacion-prototipo} de este capítulo.

\subsection{Evaluación del prototipo con usuarios}
\label{sec:evaluacion-prototipo}

La evaluación del prototipo se realizó con un grupo de 12 estudiantes matriculados en la asignatura FLP durante el período de desarrollo del trabajo de grado, quienes interactuaron con el prototipo definiendo una gramática, visualizando el AST resultante y ejecutando programas en modo paso a paso, antes de responder un instrumento compuesto por el cuestionario System Usability Scale (SUS) y un cuestionario de utilidad pedagógica de elaboración propia. Adicionalmente, se recogió la valoración del docente responsable de la asignatura mediante una entrevista semiestructurada.

\paragraph{Usabilidad (SUS).} La Tabla~\ref{tab:sus-resultados} resume los puntajes obtenidos mediante el cuestionario System Usability Scale (SUS) aplicado a los 12 estudiantes.

\begin{table}[H]
\centering
\footnotesize
\caption{Resultados del cuestionario SUS ($n = 12$)}
\label{tab:sus-resultados}
\begin{tabular}{|l|c|}
\hline
\textbf{Estadístico} & \textbf{Valor} \\
\hline
Puntaje promedio & 82.7 \\
\hline
Desviación estándar & 7.9 \\
\hline
Puntaje mínimo & 70.0 \\
\hline
Puntaje máximo & 95.0 \\
\hline
Interpretación (Bangor et al., 2009) & Excelente \\
\hline
\end{tabular}
\end{table}

El puntaje promedio obtenido (82.7 sobre 100) se ubica en la categoría ``Excelente'' de la escala de adjetivos propuesta por Bangor, Kortum y Miller, y por encima del umbral de referencia de 68 puntos que se considera aceptable para sistemas de uso general. Este resultado sugiere que la complejidad adicional que introduce el prototipo respecto a los recursos actuales de la asignatura no representa una barrera de usabilidad relevante para los estudiantes.

\paragraph{Utilidad pedagógica.} La Tabla~\ref{tab:utilidad-pedagogica} presenta la distribución de respuestas para los siete ítems del cuestionario de utilidad pedagógica, diseñado a partir de los cinco indicadores de logro seleccionados en el Capítulo~\ref{cap3}.

\begin{table}[H]
\centering
\footnotesize
\setlength{\tabcolsep}{3pt}
\caption{Distribución de respuestas - cuestionario de utilidad pedagógica ($n = 12$)}
\label{tab:utilidad-pedagogica}
\begin{tabular}{>{\raggedright\arraybackslash}p{5.2cm} >{\centering\arraybackslash}p{1.1cm} >{\centering\arraybackslash}p{1.1cm} c >{\centering\arraybackslash}p{1.1cm} >{\centering\arraybackslash}p{1.1cm} c}
\hline
\textbf{Ítem} & \textbf{Tot. desac.} & \textbf{En desac.} & \textbf{Neutral} & \textbf{De acuerdo} & \textbf{Tot. acuerdo} & \textbf{\% Acuerdo} \\
\hline
P1. Visualización del AST y comprensión de la gramática BNF & 0.0\,\% & 0.0\,\% & 8.3\,\% & 41.7\,\% & 50.0\,\% & \textbf{91.7\,\%} \\
P2. Visualización del ambiente y comprensión del estado & 0.0\,\% & 0.0\,\% & 16.7\,\% & 50.0\,\% & 33.3\,\% & \textbf{83.3\,\%} \\
P3. Claridad en la distinción ligadura/asignación & 0.0\,\% & 8.3\,\% & 25.0\,\% & 50.0\,\% & 16.7\,\% & \textbf{66.7\,\%} \\
P4. Utilidad frente a usar solo Racket y EOPL & 0.0\,\% & 0.0\,\% & 8.3\,\% & 33.3\,\% & 58.3\,\% & \textbf{91.7\,\%} \\
P5. Claridad de la ejecución paso a paso & 0.0\,\% & 0.0\,\% & 0.0\,\% & 41.7\,\% & 58.3\,\% & \textbf{100.0\,\%} \\
P6. Utilidad de la biblioteca de ejemplos & 0.0\,\% & 8.3\,\% & 16.7\,\% & 41.7\,\% & 33.3\,\% & \textbf{75.0\,\%} \\
P7. Recomendaría el prototipo a futuros estudiantes & 0.0\,\% & 0.0\,\% & 8.3\,\% & 25.0\,\% & 66.7\,\% & \textbf{91.7\,\%} \\
\hline
\end{tabular}
\end{table}

Los resultados muestran una percepción positiva y consistente en los ítems relacionados con las funcionalidades núcleo del prototipo: la visualización del AST (P1), la ejecución paso a paso (P5) y la comparación frente al uso exclusivo de Racket y EOPL (P4) concentran los niveles de acuerdo más altos, y el 91.7\,\% de los estudiantes recomendaría el prototipo a futuros cursantes de la asignatura (P7). El ítem con menor nivel de acuerdo corresponde a la distinción entre ligadura y asignación de variables (P3, 66.7\,\%), lo que es consistente con la naturaleza conceptualmente sutil de esta distinción señalada en el Capítulo~\ref{cap3}, y sugiere que la visualización del ambiente, aunque útil, no es por sí sola suficiente para resolver por completo esta dificultad específica.

En las preguntas abiertas del instrumento, la mayoría de los estudiantes destacó la posibilidad de observar el ambiente de ejecución y el AST de forma simultánea como el aporte más valioso del prototipo frente a la experiencia previa con Racket y EOPL, mientras que un grupo menor sugirió incorporar retroalimentación adicional sobre errores de sintaxis en la definición de la gramática. En la entrevista con el docente responsable se valoró positivamente la pertinencia de los conceptos representados y su viabilidad como recurso de apoyo dentro de las sesiones de clase.

\subsection{Aportes del trabajo}

Frente a los vacíos identificados en la revisión de antecedentes (Capítulo~\ref{cap2}), este trabajo aporta un prototipo que permite al estudiante definir su propia gramática en notación BNF en lugar de operar sobre un lenguaje predefinido, e integra en un mismo entorno la edición de la gramática, la visualización del AST y la inspección del ambiente de ejecución, un flujo que no se encontró combinado en ninguna de las herramientas revisadas. Adicionalmente, el trabajo aporta evidencia empírica sobre el uso de una herramienta de este tipo en un curso de lenguajes de programación de un programa de ingeniería en una universidad latinoamericana, un contexto poco representado en la literatura identificada.

\subsection{Limitaciones}

La evaluación del prototipo presenta limitaciones que deben considerarse al interpretar los resultados anteriores. El tamaño de la muestra (12 estudiantes de una única cohorte) no permite generalizar los hallazgos a otras cohortes o instituciones. La evaluación se restringió, además, a la usabilidad y a la utilidad pedagógica percibida, dejando por fuera mediciones del impacto del prototipo en el desempeño académico o en la retención del aprendizaje a largo plazo, tal como se estableció en el alcance del trabajo. Finalmente, el prototipo continúa dependiendo de Racket y EOPL como sustrato de ejecución, por lo que la carga cognitiva asociada a la sintaxis de Racket señalada en el Capítulo~\ref{cap3} no desaparece, sino que se busca mitigar mediante la representación visual complementaria.

\subsection{Trabajo futuro}

A partir de los resultados y las limitaciones identificadas, se plantean varias líneas de trabajo futuro orientadas a robustecer la evidencia empírica sobre el prototipo y a extender su alcance pedagógico. Estas líneas se desarrollan en la Sección~\ref{sec:trabajos-futuros}.
